\newpage
% \section*{ЗАКЛЮЧЕНИЕ}
\begin{center}
  \textbf{ЗАКЛЮЧЕНИЕ}
\end{center}
\addcontentsline{toc}{section}{ЗАКЛЮЧЕНИЕ}

В начале семестра, на вводной консультации, нам говорили, что этот этап научной работы обычно самый простой.
Надеюсь, что это было намеренное искажение действительности, выполненное с целью повышения нашей мотивации и веры в себя.
Другими словами -- это был обман.
Просто не было, но результат того стоил!

В рамках 2 этапа НИР была сформулирована методика оценки волнового воздействия на берег, реализация которой позволяет выделить участки береговой линии наиболее подверженные риску разрушения.
В качестве показателя воздействия волн на берег был введен CWEF (Coastal Wave Energy Flux), определяемый как проекция можности волны на нормаль к берееговой линии.

Для вычисления этого показателя в работе были рассмотренны существующие источники открытых, глобальных, геопространственных, метеорологических и батиметрическихглобальных данных.
Также были описаны инструменты доступа к ним, реализованные в проекте ссылка на который представлена на рисунке~\ref{pic:BreachTheBeach_qr_code}.
В рамках этого же проета были реализованны описанные в разделе~\ref{part:alghoritm_intro} рассчета показателя CWEF.

На основании CWEF в работе была предложена схема ранжирования, основанная на четырех взаимодополняющих критериях, каждый из которых описывает отдельный аспект волновой нагрузки: 
\begin{itemize}
    \item средний уровень энергии;
    \item экстремальную (штормовую) составляющую;
    \item направленную концентрацию волновой экспозиции;
    \item межгодовую изменчивость.
\end{itemize}

Результаты выполненных расчетов показали их непротиворечивость наблюдаемой реальности и существующей системе расселения в районе города Новороссийска и его окрестностей.
Тем не менее предложенные в разделах~\ref{sec:wer_calculation} и~\ref{sec:wer_usage} способы обощения этих показателей и способы использования являются, на данном этапе выполнения работы, лишь гипотетическими ориентирами.

Определение реальных границ для критериев риска по предложенному показателю CWEF требует более узкой спецификации задач использования морских прибрежных территорий и ретраспективного регрессионного анализа для их уточнения.

Тем не менее даже текущие результаты вселяют оптимизм и уверенность, что работа движется в праавильном направлении.
Таким образом все поставленные, в рамках 2 этипа НИР, задачи были решены в полном объеме.
